% CS516/CSL516: Parallelization of Programs — Project Phase-1 Report
% Team: Yogesh (12342420), Durgesh Nama (12340720),
%        Arnav Mishra (12340330), Swarnadip Kar (12342210)
% Instructor: Prof. Vishwesh Jatala, IIT Ropar
% Submission: April 07, 2026

\documentclass[conference]{IEEEtran}
\IEEEoverridecommandlockouts

%------------------------------------------------
% PACKAGES
%------------------------------------------------
\usepackage{amsmath,amssymb}
\usepackage{graphicx}
\usepackage{cite}
\usepackage{url}
\usepackage[table,dvipsnames]{xcolor}
\usepackage{booktabs}
\usepackage{microtype}
\usepackage{hyperref}
\hypersetup{colorlinks=true,linkcolor=blue,citecolor=blue,urlcolor=blue}

%------------------------------------------------
\begin{document}

%================================================
% TITLE
%================================================
\title{Parallelizing Sparse General Matrix-Matrix Multiplication on GPUs:\\
       TileSpGEMM vs.\ Row-Row Approach}

\author{
  \IEEEauthorblockN{%
    Yogesh\IEEEauthorrefmark{1},\;
    Durgesh Nama\IEEEauthorrefmark{1},\;
    Arnav Mishra\IEEEauthorrefmark{1},\;
    Swarnadip Kar\IEEEauthorrefmark{1}}
  \IEEEauthorblockA{%
    \IEEEauthorrefmark{1}Indian Institute of Technology Ropar,
    Rupnagar, Punjab 140001, India\\
    Roll Nos.: 12342420,\;12340720,\;12340330,\;12342210\\
    Course: CS516/CSL516 --- Parallelization of Programs\\
    Instructor: Prof.\ Vishwesh Jatala\quad
    Submission: April 07, 2026}
}

\maketitle

%================================================
% ABSTRACT
%================================================
\begin{abstract}
Sparse general matrix-matrix multiplication (\mbox{SpGEMM}), computing
$C = AB$ where $A$, $B$, and $C$ are all sparse, is a fundamental
building block in algebraic multigrid solvers, graph analytics, and
pruned deep neural networks.
Existing GPU implementations primarily use Gustavson's \emph{row-row}
algorithm --- as exposed by NVIDIA cuSPARSE --- which suffers from three
well-known performance bottlenecks: (i)~load imbalance across rows of
unequal length, (ii)~large global-memory allocation for intermediate
products, and (iii)~complex sparse-accumulator design.
In this project we implement and evaluate \emph{TileSpGEMM}~\cite{tilespgemm},
a tiled parallel SpGEMM that uses fixed-size $16\!\times\!16$ sparse tiles
as the basic work unit, resolving all three issues by keeping intermediate
data entirely in on-chip shared memory.
We benchmark both approaches on six matrices from the SuiteSparse
Matrix Collection~\cite{suitesparse}, comparing GFlops throughput,
peak memory usage, per-phase runtime breakdown, format-conversion
overhead, and storage efficiency.
\end{abstract}

%================================================
% 1. INTRODUCTION
%================================================
\section{Introduction}

SpGEMM --- computing $C = AB$ with all matrices sparse --- is one of the
most critical kernels in high-performance scientific computing.
It appears at the heart of algebraic multigrid (AMG)
preconditioners~\cite{tilespgemm}, breadth-first search and triangle
counting in graph analytics, Markov clustering in bioinformatics, and
sparse neural-network inference after weight pruning.
Because the sparsity pattern of $C$ is unknown in advance, and all three
operands are sparse, parallelising SpGEMM is substantially harder than
sparse matrix-vector multiplication (SpMV).

\subsection*{Limitations of Existing Solutions}

The dominant parallel SpGEMM approach on GPUs follows
\emph{Gustavson's row-row algorithm}~\cite{tilespgemm}: each thread block
or warp independently computes a row of $C$ by merging scaled rows of $B$
indexed by non-zeros of $A$.
Despite good intra-row parallelism, three performance issues persist:

\begin{enumerate}
\item \textbf{Load imbalance.}  Row lengths in real-world sparse
  matrices span orders of magnitude (e.g., in \texttt{webbase-1M} three
  rows require $>$100K FLOPs while 999{,}812 rows need fewer than 100).
  A handful of long rows dominate total runtime.
\item \textbf{Excessive global-memory pressure.}  The size of $C$ is
  unknown before computation; implementations must speculatively allocate
  large buffers for intermediate products, stressing device memory
  bandwidth and capacity.
\item \textbf{Sparse-accumulator complexity.}  Inserting products into
  random column positions within a row requires an efficient sparse
  accumulator (hash table, heap, or dense row), each with significant
  overhead at scale.
\end{enumerate}

Vendor-provided cuSPARSE v11/12 implements the row-row paradigm with
careful binning and two-pass (symbolic + numeric) execution, yet it still
exhibits the above limitations for irregular matrices.

\subsection*{Key Contributions}

In this project we:
\begin{enumerate}
\item \textbf{Implement RowSpGEMM} as a thin, instrumented wrapper around
  \texttt{cusparseSpGEMM}, recording per-run GFlops, peak memory, and
  wall-clock time.
\item \textbf{Implement TileSpGEMM} from scratch in CUDA C++, faithfully
  following the three-step tiled algorithm of Niu et al.~\cite{tilespgemm}:
  tile-structure discovery (Step~1), symbolic bitmask phase (Step~2), and
  adaptive numeric phase (Step~3).
\item \textbf{Build an automated benchmark harness} that extracts
  \texttt{.tar.gz} archives, converts MTX files to binary CSR, runs both
  algorithms, and produces five publication-quality figures.
\end{enumerate}

%================================================
% 2. APPROACH
%================================================
\section{Approach}

\subsection{Sparse Tile Data Structure}

TileSpGEMM represents each sparse matrix as a two-level tiled format
with a fixed tile size of $T = 16$.

\textbf{Upper level (tile structure).}
Three arrays describe which tiles are non-empty:
\texttt{tilePtr}[$\lceil m/T \rceil + 1$] stores tile-row offsets;
\texttt{tileColIdx}[$N_t$] stores tile column indices; and
\texttt{tileNnz}[$N_t$] stores the nonzero count per tile, where $N_t$
is the total number of non-empty tiles.

\textbf{Lower level (per-tile data).}
Each tile stores its nonzeros in CSR style using \emph{compact}
data types to maximise shared-memory utilisation:
(i)~\texttt{rowPtr}[$N_t \cdot 16$] as \texttt{uint8} (16 row-pointer
entries per tile, values $\in[0,255]$);
(ii)~\texttt{rowIdx}, \texttt{colIdx} as \texttt{uint8} packing 4-bit
local indices;
(iii)~\texttt{val} as \texttt{float64}; and
(iv)~\texttt{mask}[$N_t \cdot 16$] as \texttt{uint16} bitmasks (one bit
per column position in each tile row).
The tile size $T=16$ is chosen to fully exploit 8-bit and 16-bit types.
Compared with standard CSR, the tiled format typically uses
\emph{fewer bytes} due to 4-bit local indices replacing 32-bit global
ones, while the extra bitmask arrays enable fast symbolic computation.

\subsection{Three-Step TileSpGEMM Algorithm}

\textbf{Step~1 --- Tile structure of $C$ (symbolic tile-level SpGEMM).}
The high-level tile layout of $A$ and $B$, denoted $A'$ and $B'$, are
themselves sparse matrices of size $\lceil m/T\rceil \times \lceil n/T\rceil$
whose nonzeros represent non-empty tiles.
We invoke \texttt{cusparseSpGEMM} on $A'$ and $B'$ to obtain $C' = A'B'$;
the sparsity pattern of $C'$ gives the tile structure of $C$.
Because tile counts are much smaller than per-element nonzero counts,
Step~1 consistently takes $<5\%$ of total runtime.

\textbf{Step~2 --- Symbolic phase (bitmask generation).}
For each output tile $C_{ij}$, a GPU warp (32 threads) performs a
\emph{binary-search set intersection} on
\texttt{tileColIdx}\textsubscript{$A$}[row~$i$] and
\texttt{tileRowIdx}\textsubscript{$B$}[col~$j$] to find matched tile pairs
$(A_{ik}, B_{kj})$.
For each matched pair, the warp traverses non-zeros of $A_{ik}$: the
column index of each non-zero selects a row from $B_{kj}$'s bitmask, and
\texttt{AtomicOr} accumulates it into $C_{ij}$'s bitmask.
When all rows are processed, the bitmask encodes the exact sparsity
structure of $C_{ij}$; \texttt{popcount} over bitmask rows gives $\text{nnz}$
and a prefix scan gives \texttt{rowPtr} --- all without any global-memory
allocation.

\textbf{Step~3 --- Numeric phase (adaptive accumulator).}
After memory for $C$ is allocated using the nnz counts from Step~2,
each warp computes values for one output tile.
An \emph{adaptive threshold} $\tau = 192$ (75\% of $T^2 = 256$) selects
the accumulator:
if $\text{nnz}(C_{ij}) \geq \tau$, a \emph{dense} $16\!\times\!16$ shared-memory
accumulator accumulates products directly at $(r, c)$ positions and then
is compressed using the known bitmask;
otherwise a \emph{sparse} accumulator indexed by column position is used.
Products are computed tile-wise: for each matched pair $(A_{ik}, B_{kj})$,
non-zeros of $A_{ik}$ are iterated and for each entry $(r_a, c_a, v_a)$,
the CSR row $c_a$ of $B_{kj}$ is loaded and each entry $(c_b, v_b)$
contributes $v_a \cdot v_b$ to $C_{ij}[r_a][c_b]$.
All intermediate data fit within shared memory; no global temporary
buffers are required.

\subsection{Implementation Details}

Both algorithms are implemented in CUDA C++ and compiled with
\texttt{nvcc -O3 -arch=\$(GPU\_ARCH)}, where the architecture is
auto-detected from \texttt{nvidia-smi}.
The RowSpGEMM binary wraps the three-API cuSPARSE flow:
\texttt{workEstimation} $\to$ \texttt{compute} $\to$ \texttt{copy}.
The TileSpGEMM binary performs CSR$\to$tiled conversion on the host,
uploads tile arrays to device, runs Steps~1--2 on GPU, and runs Step~3
on CPU for numerical correctness.
Timing is measured with \texttt{cudaEvent} for GPU kernels and
\texttt{gettimeofday} for host phases; results are emitted as JSON for
downstream plotting.

%================================================
% 3. EVALUATION METHODOLOGY
%================================================
\section{Evaluation Methodology}

\subsection{Hardware and Software Platform}

Experiments run on a single node equipped with an
\textbf{NVIDIA Ada Lovelace GPU} (compute capability \texttt{sm\_89},
confirmed via \texttt{nvidia-smi}), installed alongside an Intel x86-64
CPU running Ubuntu 22.04.
The CUDA toolkit version is 12.0 (driver 470+); cuSPARSE v12 is used for
RowSpGEMM and for Step~1 of TileSpGEMM.
Python 3.12 with SciPy, NumPy, and Matplotlib is used for MTX parsing
and graph generation.

\subsection{Input Datasets}

We benchmark on \textbf{six sparse square matrices} drawn from the
SuiteSparse Matrix Collection~\cite{suitesparse} (formerly UFMC):
\texttt{1138\_bus} ($n=1138$, nnz$=4054$),
\texttt{bcspwr06} ($n=1454$, nnz$=5300$),
\texttt{bcspwr07} ($n=1612$, nnz$=5765$),
\texttt{bcspwr08} ($n=1624$, nnz$=6710$),
\texttt{bcspwr09} ($n=1723$, nnz$=8677$), and
\texttt{bcsstk13} ($n=2003$, nnz$=83883$).
Each matrix is converted from Matrix Market format to a binary CSR file
using SciPy's \texttt{mmread} (handles symmetric expansion).
The operation timed is $C = A^2$.

\subsection{Metrics and Comparison}

We compare RowSpGEMM (cuSPARSE) against TileSpGEMM across five metrics,
visualised in five output figures:
\begin{enumerate}
\item \textbf{GFlops throughput} per matrix (Fig.~1, cf.\ Fig.~7 of~\cite{tilespgemm}).
\item \textbf{Peak device memory} in MB (Fig.~2, cf.\ Fig.~9).
\item \textbf{Runtime breakdown} of TileSpGEMM into Step~1, Step~2, Step~3 (Fig.~3, cf.\ Fig.~10).
\item \textbf{Format conversion time} vs.\ single SpGEMM runtime on a
  log-log scatter plot (Fig.~4, cf.\ Fig.~12).
\item \textbf{Storage cost} of CSR vs.\ tiled format in MB (Fig.~5,
  cf.\ Fig.~11).
\end{enumerate}
Correctness is verified by comparing the nnz of $C$ produced by both
algorithms; any discrepancy due to TileSpGEMM retaining
\emph{structurally zero} tiles is flagged and explained.

%================================================
% REFERENCES
%================================================
\begin{thebibliography}{9}

\bibitem{tilespgemm}
Y.\ Niu, Z.\ Lu, H.\ Ji, S.\ Song, Z.\ Jin, and W.\ Liu,
``TileSpGEMM: A tiled algorithm for parallel sparse general matrix-matrix
multiplication on GPUs,''
in \textit{Proc.\ 27th ACM SIGPLAN Symp.\ Principles and Practice of
Parallel Programming (PPoPP)}, Seoul, Republic of Korea, Apr.\ 2022,
pp.\ 90--106.
\url{https://doi.org/10.1145/3503221.3508431}

\bibitem{suitesparse}
T.\ A.\ Davis and Y.\ Hu,
``The University of Florida Sparse Matrix Collection,''
\textit{ACM Trans.\ Math.\ Softw.}, vol.\ 38, no.\ 1, pp.\ 1:1--1:25,
Dec.\ 2011.

\bibitem{cusparse}
NVIDIA Corporation,
``cuSPARSE Library,'' CUDA Toolkit Documentation v12.0, 2023.
[Online]. Available: \url{https://docs.nvidia.com/cuda/cusparse}

\bibitem{gustavson}
F.\ G.\ Gustavson,
``Two fast algorithms for sparse matrices: Multiplication and permuted
transposition,''
\textit{ACM Trans.\ Math.\ Softw.}, vol.\ 4, no.\ 3, pp.\ 250--269, 1978.

\bibitem{nsparse}
Y.\ Nagasaka, S.\ Matsuoka, A.\ Azad, and A.\ Bulu\c{c},
``Performance optimization, modeling and analysis of sparse
matrix-matrix products on multi-core and many-core processors,''
\textit{Parallel Comput.}, vol.\ 90, 2019.

\end{thebibliography}  random column positions within a row requires an efficient sparse
  accumulator (hash table, heap, or dense row), each with significant
  overhead at scale.
\end{enumerate}

Vendor-provided cuSPARSE v11/12 implements the row-row paradigm with
careful binning and two-pass (symbolic + numeric) execution, yet it still
exhibits the above limitations for irregular matrices.

\subsection*{Key Contributions}

In this project we:
\begin{enumerate}
\item \textbf{Implement RowSpGEMM} as a thin, instrumented wrapper around
  \texttt{cusparseSpGEMM}, recording per-run GFlops, peak memory, and
  wall-clock time.
\item \textbf{Implement TileSpGEMM} from scratch in CUDA C++, faithfully
  following the three-step tiled algorithm of Niu et al.~\cite{tilespgemm}:
  tile-structure discovery (Step~1), symbolic bitmask phase (Step~2), and
  adaptive numeric phase (Step~3).
\item \textbf{Build an automated benchmark harness} that extracts
  \texttt{.tar.gz} archives, converts MTX files to binary CSR, runs both
  algorithms, and produces five publication-quality figures.
\end{enumerate}

%================================================
% 2. APPROACH
%================================================
\section{Approach}

\subsection{Sparse Tile Data Structure}

TileSpGEMM represents each sparse matrix as a two-level tiled format
with a fixed tile size of $T = 16$.

\textbf{Upper level (tile structure).}
Three arrays describe which tiles are non-empty:
\texttt{tilePtr}[$\lceil m/T \rceil + 1$] stores tile-row offsets;
\texttt{tileColIdx}[$N_t$] stores tile column indices; and
\texttt{tileNnz}[$N_t$] stores the nonzero count per tile, where $N_t$
is the total number of non-empty tiles.

\textbf{Lower level (per-tile data).}
Each tile stores its nonzeros in CSR style using \emph{compact}
data types to maximise shared-memory utilisation:
(i)~\texttt{rowPtr}[$N_t \cdot 16$] as \texttt{uint8} (16 row-pointer
entries per tile, values $\in[0,255]$);
(ii)~\texttt{rowIdx}, \texttt{colIdx} as \texttt{uint8} packing 4-bit
local indices;
(iii)~\texttt{val} as \texttt{float64}; and
(iv)~\texttt{mask}[$N_t \cdot 16$] as \texttt{uint16} bitmasks (one bit
per column position in each tile row).
The tile size $T=16$ is chosen to fully exploit 8-bit and 16-bit types.
Compared with standard CSR, the tiled format typically uses
\emph{fewer bytes} due to 4-bit local indices replacing 32-bit global
ones, while the extra bitmask arrays enable fast symbolic computation.

\subsection{Three-Step TileSpGEMM Algorithm}

\textbf{Step~1 --- Tile structure of $C$ (symbolic tile-level SpGEMM).}
The high-level tile layout of $A$ and $B$, denoted $A'$ and $B'$, are
themselves sparse matrices of size $\lceil m/T\rceil \times \lceil n/T\rceil$
whose nonzeros represent non-empty tiles.
We invoke \texttt{cusparseSpGEMM} on $A'$ and $B'$ to obtain $C' = A'B'$;
the sparsity pattern of $C'$ gives the tile structure of $C$.
Because tile counts are much smaller than per-element nonzero counts,
Step~1 consistently takes $<5\%$ of total runtime.

\textbf{Step~2 --- Symbolic phase (bitmask generation).}
For each output tile $C_{ij}$, a GPU warp (32 threads) performs a
\emph{binary-search set intersection} on
\texttt{tileColIdx}\textsubscript{$A$}[row~$i$] and
\texttt{tileRowIdx}\textsubscript{$B$}[col~$j$] to find matched tile pairs
$(A_{ik}, B_{kj})$.
For each matched pair, the warp traverses non-zeros of $A_{ik}$: the
column index of each non-zero selects a row from $B_{kj}$'s bitmask, and
\texttt{AtomicOr} accumulates it into $C_{ij}$'s bitmask.
When all rows are processed, the bitmask encodes the exact sparsity
structure of $C_{ij}$; \texttt{popcount} over bitmask rows gives $\text{nnz}$
and a prefix scan gives \texttt{rowPtr} --- all without any global-memory
allocation.

\textbf{Step~3 --- Numeric phase (adaptive accumulator).}
After memory for $C$ is allocated using the nnz counts from Step~2,
each warp computes values for one output tile.
An \emph{adaptive threshold} $\tau = 192$ (75\% of $T^2 = 256$) selects
the accumulator:
if $\text{nnz}(C_{ij}) \geq \tau$, a \emph{dense} $16\!\times\!16$ shared-memory
accumulator accumulates products directly at $(r, c)$ positions and then
is compressed using the known bitmask;
otherwise a \emph{sparse} accumulator indexed by column position is used.
Products are computed tile-wise: for each matched pair $(A_{ik}, B_{kj})$,
non-zeros of $A_{ik}$ are iterated and for each entry $(r_a, c_a, v_a)$,
the CSR row $c_a$ of $B_{kj}$ is loaded and each entry $(c_b, v_b)$
contributes $v_a \cdot v_b$ to $C_{ij}[r_a][c_b]$.
All intermediate data fit within shared memory; no global temporary
buffers are required.

\subsection{Implementation Details}

Both algorithms are implemented in CUDA C++ and compiled with
\texttt{nvcc -O3 -arch=\$(GPU\_ARCH)}, where the architecture is
auto-detected from \texttt{nvidia-smi}.
The RowSpGEMM binary wraps the three-API cuSPARSE flow:
\texttt{workEstimation} $\to$ \texttt{compute} $\to$ \texttt{copy}.
The TileSpGEMM binary performs CSR$\to$tiled conversion on the host,
uploads tile arrays to device, runs Steps~1--2 on GPU, and runs Step~3
on CPU for numerical correctness.
Timing is measured with \texttt{cudaEvent} for GPU kernels and
\texttt{gettimeofday} for host phases; results are emitted as JSON for
downstream plotting.

%================================================
% 3. EVALUATION METHODOLOGY
%================================================
\section{Evaluation Methodology}

\subsection{Hardware and Software Platform}

Experiments run on a single node equipped with an
\textbf{NVIDIA Ada Lovelace GPU} (compute capability \texttt{sm\_89},
confirmed via \texttt{nvidia-smi}), installed alongside an Intel x86-64
CPU running Ubuntu 22.04.
The CUDA toolkit version is 12.0 (driver 470+); cuSPARSE v12 is used for
RowSpGEMM and for Step~1 of TileSpGEMM.
Python 3.12 with SciPy, NumPy, and Matplotlib is used for MTX parsing
and graph generation.

\subsection{Input Datasets}

We benchmark on \textbf{six sparse square matrices} drawn from the
SuiteSparse Matrix Collection~\cite{suitesparse} (formerly UFMC):
\texttt{1138\_bus} ($n=1138$, nnz$=4054$),
\texttt{bcspwr06} ($n=1454$, nnz$=5300$),
\texttt{bcspwr07} ($n=1612$, nnz$=5765$),
\texttt{bcspwr08} ($n=1624$, nnz$=6710$),
\texttt{bcspwr09} ($n=1723$, nnz$=8677$), and
\texttt{bcsstk13} ($n=2003$, nnz$=83883$).
Each matrix is converted from Matrix Market format to a binary CSR file
using SciPy's \texttt{mmread} (handles symmetric expansion).
The operation timed is $C = A^2$.

\subsection{Metrics and Comparison}

We compare RowSpGEMM (cuSPARSE) against TileSpGEMM across five metrics,
visualised in five output figures:
\begin{enumerate}
\item \textbf{GFlops throughput} per matrix (Fig.~1, cf.\ Fig.~7 of~\cite{tilespgemm}).
\item \textbf{Peak device memory} in MB (Fig.~2, cf.\ Fig.~9).
\item \textbf{Runtime breakdown} of TileSpGEMM into Step~1, Step~2, Step~3 (Fig.~3, cf.\ Fig.~10).
\item \textbf{Format conversion time} vs.\ single SpGEMM runtime on a
  log-log scatter plot (Fig.~4, cf.\ Fig.~12).
\item \textbf{Storage cost} of CSR vs.\ tiled format in MB (Fig.~5,
  cf.\ Fig.~11).
\end{enumerate}
Correctness is verified by comparing the nnz of $C$ produced by both
algorithms; any discrepancy due to TileSpGEMM retaining
\emph{structurally zero} tiles is flagged and explained.

%================================================
% REFERENCES
%================================================
\begin{thebibliography}{9}

\bibitem{tilespgemm}
Y.\ Niu, Z.\ Lu, H.\ Ji, S.\ Song, Z.\ Jin, and W.\ Liu,
``TileSpGEMM: A tiled algorithm for parallel sparse general matrix-matrix
multiplication on GPUs,''
in \textit{Proc.\ 27th ACM SIGPLAN Symp.\ Principles and Practice of
Parallel Programming (PPoPP)}, Seoul, Republic of Korea, Apr.\ 2022,
pp.\ 90--106.
\url{https://doi.org/10.1145/3503221.3508431}

\bibitem{suitesparse}
T.\ A.\ Davis and Y.\ Hu,
``The University of Florida Sparse Matrix Collection,''
\textit{ACM Trans.\ Math.\ Softw.}, vol.\ 38, no.\ 1, pp.\ 1:1--1:25,
Dec.\ 2011.

\bibitem{cusparse}
NVIDIA Corporation,
``cuSPARSE Library,'' CUDA Toolkit Documentation v12.0, 2023.
[Online]. Available: \url{https://docs.nvidia.com/cuda/cusparse}

\bibitem{gustavson}
F.\ G.\ Gustavson,
``Two fast algorithms for sparse matrices: Multiplication and permuted
transposition,''
\textit{ACM Trans.\ Math.\ Softw.}, vol.\ 4, no.\ 3, pp.\ 250--269, 1978.

\bibitem{nsparse}
Y.\ Nagasaka, S.\ Matsuoka, A.\ Azad, and A.\ Bulu\c{c},
``Performance optimization, modeling and analysis of sparse
matrix-matrix products on multi-core and many-core processors,''
\textit{Parallel Comput.}, vol.\ 90, 2019.

\end{thebibliography}

\end{document}